% Appendix H -- Latency (body). Sources: rss_ext/DESIGN_ext.md 1.3-1.5, rss_ext/E8_latency/REPORT.md, E8c_latency/REPORT.md,
%   main/reports/E7_tdmpc2.md, posthoc_tdmpc2.md, tmlr_ext/K_kernel/REPORT.md 2. Tables: tables/tabH_e8.tex, tabH_e17.tex.
\textbf{Protocol.} All latencies are of the TD-MPC2 world model, measured on one GPU used exclusively by a single process. After 30 warm-up calls, each condition is timed in blocks with CUDA events; a block lasts about two seconds or more (at least 10 calls) and its value is the median time per call. A rounding variant and \dtm are timed as two adjacent blocks in random order, which form a pair; the float formats are single blocks. Ten rounds are run, with the order of the conditions reshuffled in every round, and the SM clock, GPU temperature and host load are recorded for every block. A condition's latency is the median of its block values, a ratio is a ratio of these medians, and an overhead is the median over the ten pairs of (variant $-$ \dtm)/\dtm, with a bootstrap interval over pairs ($10^4$ resamples). ``Graph'' captures the world-model rollouts of one control step (6 iterations $\times$ 3 horizon steps of dynamics and reward for $N$ candidates, plus the policy rollouts) as one CUDA graph, as TD-MPC2's compiled planner runs them; ``step'' captures the whole planning step, with encoder, policy, Q-ensemble and MPPI update, in one graph (E8c). E8 and E8c time TD-MPC2's architecture at 5M, 19M, 48M and 317M parameters with random weights, since latency does not depend on the weight values.

\textbf{Eager execution and two measurement modes.} E7 timed the fused Triton kernel eagerly on an RTX A5000 at $N=512$: the integer kernel was \nv{b.h5.slow}\% slower than FP32 (median over conditions, per kernel), and the dither overhead lay between \nv{b.h5.min} and \nv{b.h5.max}\% (H5 held, with an upper limit of 10\%). These eager timings had two modes: in a post hoc paired re-measurement, \nv{h.ph.slow}\% of the blocks were at least 15\% slower than the rest, about half of them together with a drop of the SM clock, and the paired overhead medians lay within \nv{h.ph.med}\%, with intervals containing zero in \nv{h.ph.ci0} of \nv{h.ph.n} conditions. Eager execution is dominated by kernel-launch cost at small sizes: E8 found the eager integer rollout at 5M and $N=512$ at \nv{h.e8.eager5}$\times$ the FP32 time (\nv{h.e8.eager.range}$\times$ over all sizes and $N$), and the eager planning step at 5M \nv{h.e8.step5}\% slower than FP32 with the new kernel (\nv{h.e8.step5.e7}\% with E7's), while at 317M the eager step took \nv{h.e8.step317}$\times$ the FP32 time and \nv{h.e8.step317.min}$\times$ that of the faster of TF32 and BF16. Captured as CUDA graphs, the integer path is faster than FP32 in every condition (Table~\ref{tab:e8}).

\begin{table}[h]
\caption{World-model latency on an NVIDIA A40 as CUDA graphs: E8 rollouts (FP32 time, integer W8A4 time relative to FP32 and to the faster of TF32 and BF16 with per-call conversion, \shiter overhead) and E8c, a second run on another A40 (integer rollout relative to BF16 with weights converted once in advance; whole planning step relative to FP32 and to pre-converted BF16, the fastest float format there, and \shiter overhead).}
\label{tab:e8}
\centering\scriptsize
\setlength{\tabcolsep}{4pt}
\begin{tabular}{@{}lrrrrrrrrr@{}}
\toprule
 & & \multicolumn{4}{c}{E8 rollout} & E8c rollout & \multicolumn{3}{c}{E8c planning step} \\
\cmidrule(lr){3-6}\cmidrule(lr){7-7}\cmidrule(l){8-10}
Size & $N$ & FP32 (ms) & /FP32 & /min(TF32, BF16) & \shiter (\%) & /BF16 pre-conv. & /FP32 & /BF16 pre-conv. & \shiter (\%) \\
\midrule
\inputrows{tables/tabH_e8.tex}
\bottomrule
\end{tabular}
\end{table}

\textbf{The first layer.} The integer path keeps the first layer in FP32, and much of its gap to pre-converted BF16 lies there: with TF32 in that layer (post hoc, not evaluated in closed loop), the integer rollout takes \nv{e8c.tf32first.over.bf16}$\times$ the time of pre-converted BF16.

\textbf{Integer \detshift and the cost of the offset (E17).} Table~\ref{tab:latency} gives the integer-kernel overheads of E17 (paired blocks, 10 rounds per condition, median and 95\% interval). The variant \shiter adds one draw per layer and call that every candidate shares, so its overhead is at most \nv{e17.shiter.overhead.max}\%. The variant \detshift adds the same draw and, in addition, the correction $W(\Delta o)$ of the next layer's bias, a matrix--vector product that reads that layer's int8 weights once more per call. This product does not grow with $N$, so its share is largest at $N=512$, where \detshift adds \nv{e17.detshift.5m512}\% (5M) and \nv{e17.detshift.317m512}\% (317M), against \nv{e17.detshift.big}\% at $N\ge2048$. The split of this product across GPU programs was changed once after the smoke test, before the main measurement.

\begin{table}[h]
\caption{World-model rollout latency of the integer kernel (\dtm) and the overheads of \shiter and \detshift (E17, NVIDIA A40, CUDA graph).}
\label{tab:latency}
\centering\small
\begin{tabular}{@{}lrrrr@{}}
\toprule
Model & $N$ & \dtm (ms) & \shiter overhead (\%) & \detshift overhead (\%) \\
\midrule
\inputrows{tables/tabH_e17.tex}
\bottomrule
\end{tabular}
\end{table}

\textbf{Hardware with fast low-bit arithmetic.} On the A40, an Ampere GPU, int8 tensor cores do not beat BF16 with pre-converted weights in our rollouts. On an RTX 5090 (Blackwell), E8b timed single GEMMs at the 317M layer width with the PyTorch of a CUDA 12.8 image: NVFP4 was up to \nv{e8b.nvfp4.gemm}$\times$ faster than BF16. We did not run closed-loop planning on that GPU.
\clearpage
